\documentclass[10pt,a4paper]{amsart}
\usepackage{amsmath,amssymb,amsthm,graphicx}
\usepackage{verbatim}
\newenvironment{reference}{\comment}{\endcomment}
%\newenvironment{reference}{}{}
\newenvironment{slogan}{\comment}{\endcomment}
\newenvironment{history}{\comment}{\endcomment}

% For commutative diagrams we use Xy-pic
\usepackage[all]{xy}

% We use 2cell for 2-commutative diagrams.
\xyoption{2cell}
% Theorem environments.
%
\theoremstyle{plain}
\newtheorem{theorem}[subsection]{Theorem}
\newtheorem{proposition}[subsection]{Proposition}
\newtheorem{lemma}[subsection]{Lemma}

\theoremstyle{definition}
\newtheorem{definition}[subsection]{Definition}
\newtheorem{example}[subsection]{Example}
\newtheorem{exercise}[subsection]{Exercise}
\newtheorem{situation}[subsection]{Situation}

\theoremstyle{remark}
\newtheorem{remark}[subsection]{Remark}
\newtheorem{remarks}[subsection]{Remarks}

\numberwithin{equation}{subsection}

% Macros
%
\def\lim{\mathop{\mathrm{lim}}\nolimits}
\def\colim{\mathop{\mathrm{colim}}\nolimits}
\def\Spec{\mathop{\mathrm{Spec}}}
\def\Hom{\mathop{\mathrm{Hom}}\nolimits}
\def\Ext{\mathop{\mathrm{Ext}}\nolimits}
\def\SheafHom{\mathop{\mathcal{H}\!\mathit{om}}\nolimits}
\def\SheafExt{\mathop{\mathcal{E}\!\mathit{xt}}\nolimits}
\def\Sch{\mathit{Sch}}
\def\Mor{\mathop{\mathrm{Mor}}\nolimits}
\def\Ob{\mathop{\mathrm{Ob}}\nolimits}
\def\Sh{\mathop{\mathit{Sh}}\nolimits}
\def\NL{\mathop{N\!L}\nolimits}
\def\CH{\mathop{\mathrm{CH}}\nolimits}
\def\proetale{{pro\text{-}\acute{e}tale}}
\def\etale{{\acute{e}tale}}
\def\QCoh{\mathit{QCoh}}
\def\Ker{\mathop{\mathrm{Ker}}}
\def\Im{\mathop{\mathrm{Im}}}
\def\Coker{\mathop{\mathrm{Coker}}}
\def\Coim{\mathop{\mathrm{Coim}}}

% Boxtimes
%
\DeclareMathSymbol{\boxtimes}{\mathbin}{AMSa}{"02}

%
% Macros for moduli stacks/spaces
%
\def\QCohstack{\mathcal{QC}\!\mathit{oh}}
\def\Cohstack{\mathcal{C}\!\mathit{oh}}
\def\Spacesstack{\mathcal{S}\!\mathit{paces}}
\def\Quotfunctor{\mathrm{Quot}}
\def\Hilbfunctor{\mathrm{Hilb}}
\def\Curvesstack{\mathcal{C}\!\mathit{urves}}
\def\Polarizedstack{\mathcal{P}\!\mathit{olarized}}
\def\Complexesstack{\mathcal{C}\!\mathit{omplexes}}
% \Pic is the operator that assigns to X its picard group, usage \Pic(X)
% \Picardstack_{X/B} denotes the Picard stack of X over B
% \Picardfunctor_{X/B} denotes the Picard functor of X over B
\def\Pic{\mathop{\mathrm{Pic}}\nolimits}
\def\Picardstack{\mathcal{P}\!\mathit{ic}}
\def\Picardfunctor{\mathrm{Pic}}
\def\Deformationcategory{\mathcal{D}\!\mathit{ef}}
\usepackage{cleveref}
\crefname{originallemma}{original lemma}{original lemmas}
\Crefname{originallemma}{Original lemma}{Original lemmas}
\makeatletter
\newlabel{algebra-lemma-apply-grothendieck}{{0337}{source}{}{sourcepointer-1}{}}
\newlabel{algebra-lemma-apply-grothendieck@cref}{{[originallemma][0][]0337}{[1][0][]source}{}{}{}}
\makeatother
\crefname{originallemma}{original lemma}{original lemmas}
\Crefname{originallemma}{Original lemma}{Original lemmas}
\makeatletter
\newlabel{algebra-lemma-characterize-separable-field-extensions}{{030W}{source}{}{sourcepointer-2}{}}
\newlabel{algebra-lemma-characterize-separable-field-extensions@cref}{{[originallemma][0][]030W}{[1][0][]source}{}{}{}}
\makeatother
\crefname{originallemma}{original lemma}{original lemmas}
\Crefname{originallemma}{Original lemma}{Original lemmas}
\makeatletter
\newlabel{algebra-lemma-criterion-reduced}{{031R}{source}{}{sourcepointer-3}{}}
\newlabel{algebra-lemma-criterion-reduced@cref}{{[originallemma][0][]031R}{[1][0][]source}{}{}{}}
\makeatother
\crefname{originallemma}{original lemma}{original lemmas}
\Crefname{originallemma}{Original lemma}{Original lemmas}
\makeatletter
\newlabel{algebra-lemma-descent-reduced}{{033F}{source}{}{sourcepointer-4}{}}
\newlabel{algebra-lemma-descent-reduced@cref}{{[originallemma][0][]033F}{[1][0][]source}{}{}{}}
\makeatother
\crefname{originalsection}{original section}{original sections}
\Crefname{originalsection}{Original section}{Original sections}
\makeatletter
\newlabel{algebra-section-differentials}{{00RM}{source}{}{sourcepointer-5}{}}
\newlabel{algebra-section-differentials@cref}{{[originalsection][0][]00RM}{[1][0][]source}{}{}{}}
\makeatother
\crefname{originalsection}{original section}{original sections}
\Crefname{originalsection}{Original section}{Original sections}
\makeatletter
\newlabel{algebra-section-netherlander}{{00S0}{source}{}{sourcepointer-6}{}}
\newlabel{algebra-section-netherlander@cref}{{[originalsection][0][]00S0}{[1][0][]source}{}{}{}}
\makeatother
\crefname{originallemma}{original lemma}{original lemmas}
\Crefname{originallemma}{Original lemma}{Original lemmas}
\makeatletter
\newlabel{divisors-lemma-base-change-and-fitting-ideal-omega}{{0C3I}{source}{}{sourcepointer-7}{}}
\newlabel{divisors-lemma-base-change-and-fitting-ideal-omega@cref}{{[originallemma][0][]0C3I}{[1][0][]source}{}{}{}}
\makeatother
\crefname{originalreference}{original reference}{original references}
\Crefname{originalreference}{Original reference}{Original references}
\makeatletter
\newlabel{lemma-2-branches-delta-1}{{0C4A}{source}{}{sourcepointer-8}{}}
\newlabel{lemma-2-branches-delta-1@cref}{{[originalreference][0][]0C4A}{[1][0][]source}{}{}{}}
\makeatother
\crefname{originalreference}{original reference}{original references}
\Crefname{originalreference}{Original reference}{Original references}
\makeatletter
\newlabel{lemma-fitting-ideal}{{07Z8}{source}{}{sourcepointer-9}{}}
\newlabel{lemma-fitting-ideal@cref}{{[originalreference][0][]07Z8}{[1][0][]source}{}{}{}}
\makeatother
\crefname{originalreference}{original reference}{original references}
\Crefname{originalreference}{Original reference}{Original references}
\makeatletter
\newlabel{lemma-multicross-algebra}{{0C1V}{source}{}{sourcepointer-10}{}}
\newlabel{lemma-multicross-algebra@cref}{{[originalreference][0][]0C1V}{[1][0][]source}{}{}{}}
\makeatother
\crefname{originalreference}{original reference}{original references}
\Crefname{originalreference}{Original reference}{Original references}
\makeatletter
\newlabel{lemma-nodal-algebraic}{{0C49}{source}{}{sourcepointer-11}{}}
\newlabel{lemma-nodal-algebraic@cref}{{[originalreference][0][]0C49}{[1][0][]source}{}{}{}}
\makeatother
\crefname{originallemma}{original lemma}{original lemmas}
\Crefname{originallemma}{Original lemma}{Original lemmas}
\makeatletter
\newlabel{varieties-lemma-delta-invariant-and-change-of-fields}{{0C3X}{source}{}{sourcepointer-12}{}}
\newlabel{varieties-lemma-delta-invariant-and-change-of-fields@cref}{{[originallemma][0][]0C3X}{[1][0][]source}{}{}{}}
\makeatother
\crefname{originallemma}{original lemma}{original lemmas}
\Crefname{originallemma}{Original lemma}{Original lemmas}
\makeatletter
\newlabel{varieties-lemma-delta-invariant-and-change-of-fields-better}{{0C3Y}{source}{}{sourcepointer-13}{}}
\newlabel{varieties-lemma-delta-invariant-and-change-of-fields-better@cref}{{[originallemma][0][]0C3Y}{[1][0][]source}{}{}{}}
\makeatother
\crefname{originallemma}{original lemma}{original lemmas}
\Crefname{originallemma}{Original lemma}{Original lemmas}
\makeatletter
\newlabel{varieties-lemma-delta-invariant-is-zero}{{0C3U}{source}{}{sourcepointer-14}{}}
\newlabel{varieties-lemma-delta-invariant-is-zero@cref}{{[originallemma][0][]0C3U}{[1][0][]source}{}{}{}}
\makeatother
\crefname{originallemma}{original lemma}{original lemmas}
\Crefname{originallemma}{Original lemma}{Original lemmas}
\makeatletter
\newlabel{varieties-lemma-geometric-branches-and-change-of-fields}{{0C40}{source}{}{sourcepointer-15}{}}
\newlabel{varieties-lemma-geometric-branches-and-change-of-fields@cref}{{[originallemma][0][]0C40}{[1][0][]source}{}{}{}}
\makeatother
\crefname{originallemma}{original lemma}{original lemmas}
\Crefname{originallemma}{Original lemma}{Original lemmas}
\makeatletter
\newlabel{varieties-lemma-geometrically-normal-in-codim-1}{{0C3P}{source}{}{sourcepointer-16}{}}
\newlabel{varieties-lemma-geometrically-normal-in-codim-1@cref}{{[originallemma][0][]0C3P}{[1][0][]source}{}{}{}}
\makeatother
\crefname{originallemma}{original lemma}{original lemmas}
\Crefname{originallemma}{Original lemma}{Original lemmas}
\makeatletter
\newlabel{varieties-lemma-geometrically-reduced-upstairs}{{0384}{source}{}{sourcepointer-17}{}}
\newlabel{varieties-lemma-geometrically-reduced-upstairs@cref}{{[originallemma][0][]0384}{[1][0][]source}{}{}{}}
\makeatother
\title{Six equivalent characterizations of a node of a 1-dimensional algebraic scheme at a closed point}
\author{The Stacks Project Authors}
\date{stacks-project:curves.tex:local-snapshot-20261012; GNU FDL 1.2 or later}
\begin{document}
\maketitle
\noindent\textit{Editorial note.} All delivered mathematical text is the author's own, copied from the pinned Stacks Project chapter file curves.tex; no statement or proof step was rewritten, reordered or removed.\par\medskip
\noindent\textit{Editorial note.} The counted proof is the author's complete argument: the algebraically closed case with the quadric normal form, the Fitting-ideal presentation, the separability descent, the flat transfer of conditions (4) and (5), the reduction to the geometric fibre in both directions, and the delta-invariant inequality with its strictness argument are all retained verbatim.\par\medskip
\noindent\textit{Editorial note.} The label lemma-nodal names the counted statement, and the proof environment immediately following it in the source is delivered verbatim as its proof.\par\medskip
\noindent\textit{Editorial note.} The counted statement is carried by the author's own lemma environment in the chapter source.\par\medskip
\noindent\textit{Allowed standard prerequisites (editorial).}\begin{itemize}
\item The first Fitting ideal of a presented module. Same chapter, Lemma lemma-fitting-ideal, with the standard presentation R\textasciicircum{}m -> R\textasciicircum{}n -> Omega\_\{R/k\} -> 0 from the partial-derivative matrix. Supplies the computational core: conditions (2), (4), (5) are equivalent because the first Fitting ideal is generated by the (n-1)-minors of the Jacobian matrix.
\item Nodal algebras and the delta-invariant. Same chapter, Lemmas lemma-nodal-algebraic, lemma-multicross-algebra and lemma-2-branches-delta-1. Identifies condition (1) over an algebraically closed base with the nondegenerate quadric xy, and condition (6) with separability, reducedness, delta-invariant 1 and 2 geometric branches.
\end{itemize}\medskip
\ifdefined\hypertarget\hypertarget{sourcepointer-1}{}\fi
\noindent\textit{Source reference.} Original lemma 0337: sources/stacks-chow/tags.txt, source lines 4046--4046. Prerequisite pointer invoked at proof line(s) 4464; it names a standard fact of a sibling chapter that the counted argument uses, and every step of the counted proof built on it is retained verbatim in the delivered body.\par\medskip
\ifdefined\hypertarget\hypertarget{sourcepointer-2}{}\fi
\noindent\textit{Source reference.} Original lemma 030W: sources/stacks-chow/tags.txt, source lines 3963--3963. Prerequisite pointer invoked at proof line(s) 4455; it names a standard fact of a sibling chapter that the counted argument uses, and every step of the counted proof built on it is retained verbatim in the delivered body.\par\medskip
\ifdefined\hypertarget\hypertarget{sourcepointer-3}{}\fi
\noindent\textit{Source reference.} Original lemma 031R: sources/stacks-chow/tags.txt, source lines 3994--3994. Prerequisite pointer invoked at proof line(s) 4466; it names a standard fact of a sibling chapter that the counted argument uses, and every step of the counted proof built on it is retained verbatim in the delivered body.\par\medskip
\ifdefined\hypertarget\hypertarget{sourcepointer-4}{}\fi
\noindent\textit{Source reference.} Original lemma 033F: sources/stacks-chow/tags.txt, source lines 4054--4054. Prerequisite pointer invoked at proof line(s) 4463; it names a standard fact of a sibling chapter that the counted argument uses, and every step of the counted proof built on it is retained verbatim in the delivered body.\par\medskip
\ifdefined\hypertarget\hypertarget{sourcepointer-5}{}\fi
\noindent\textit{Source reference.} Original section 00RM: sources/stacks-chow/tags.txt, source lines 1008--1008. Prerequisite pointer invoked at proof line(s) 4427; it names a standard fact of a sibling chapter that the counted argument uses, and every step of the counted proof built on it is retained verbatim in the delivered body.\par\medskip
\ifdefined\hypertarget\hypertarget{sourcepointer-6}{}\fi
\noindent\textit{Source reference.} Original section 00S0: sources/stacks-chow/tags.txt, source lines 1022--1022. Prerequisite pointer invoked at proof line(s) 4428; it names a standard fact of a sibling chapter that the counted argument uses, and every step of the counted proof built on it is retained verbatim in the delivered body.\par\medskip
\ifdefined\hypertarget\hypertarget{sourcepointer-7}{}\fi
\noindent\textit{Source reference.} Original lemma 0C3I: sources/stacks-chow/tags.txt, source lines 15193--15193. Prerequisite pointer invoked at proof line(s) 4443; it names a standard fact of a sibling chapter that the counted argument uses, and every step of the counted proof built on it is retained verbatim in the delivered body.\par\medskip
\ifdefined\hypertarget\hypertarget{sourcepointer-8}{}\fi
\noindent\textit{Source reference.} Original reference 0C4A: sources/stacks-curves/pinned-original-source.tex, source lines 4181--4181. Prerequisite pointer invoked at proof line(s) 4484; it names a standard fact of the same chapter that the counted argument uses, and every step of the counted proof built on it is retained verbatim in the delivered body.\par\medskip
\ifdefined\hypertarget\hypertarget{sourcepointer-9}{}\fi
\noindent\textit{Source reference.} Original reference 07Z8: sources/stacks-curves/pinned-original-source.tex, source lines 4312--4312. Prerequisite pointer invoked at proof line(s) 4432; it names a standard fact of the same chapter that the counted argument uses, and every step of the counted proof built on it is retained verbatim in the delivered body.\par\medskip
\ifdefined\hypertarget\hypertarget{sourcepointer-10}{}\fi
\noindent\textit{Source reference.} Original reference 0C1V: sources/stacks-curves/pinned-original-source.tex, source lines 3431--3431. Prerequisite pointer invoked at proof line(s) 4416; it names a standard fact of the same chapter that the counted argument uses, and every step of the counted proof built on it is retained verbatim in the delivered body.\par\medskip
\ifdefined\hypertarget\hypertarget{sourcepointer-11}{}\fi
\noindent\textit{Source reference.} Original reference 0C49: sources/stacks-curves/pinned-original-source.tex, source lines 4090--4090. Prerequisite pointer invoked at proof line(s) 4396; it names a standard fact of the same chapter that the counted argument uses, and every step of the counted proof built on it is retained verbatim in the delivered body.\par\medskip
\ifdefined\hypertarget\hypertarget{sourcepointer-12}{}\fi
\noindent\textit{Source reference.} Original lemma 0C3X: sources/stacks-chow/tags.txt, source lines 15207--15207. Prerequisite pointer invoked at proof line(s) 4501; it names a standard fact of a sibling chapter that the counted argument uses, and every step of the counted proof built on it is retained verbatim in the delivered body.\par\medskip
\ifdefined\hypertarget\hypertarget{sourcepointer-13}{}\fi
\noindent\textit{Source reference.} Original lemma 0C3Y: sources/stacks-chow/tags.txt, source lines 15208--15208. Prerequisite pointer invoked at proof line(s) 4513; it names a standard fact of a sibling chapter that the counted argument uses, and every step of the counted proof built on it is retained verbatim in the delivered body.\par\medskip
\ifdefined\hypertarget\hypertarget{sourcepointer-14}{}\fi
\noindent\textit{Source reference.} Original lemma 0C3U: sources/stacks-chow/tags.txt, source lines 15204--15204. Prerequisite pointer invoked at proof line(s) 4506; it names a standard fact of a sibling chapter that the counted argument uses, and every step of the counted proof built on it is retained verbatim in the delivered body.\par\medskip
\ifdefined\hypertarget\hypertarget{sourcepointer-15}{}\fi
\noindent\textit{Source reference.} Original lemma 0C40: sources/stacks-chow/tags.txt, source lines 15210--15210. Prerequisite pointer invoked at proof line(s) 4492; it names a standard fact of a sibling chapter that the counted argument uses, and every step of the counted proof built on it is retained verbatim in the delivered body.\par\medskip
\ifdefined\hypertarget\hypertarget{sourcepointer-16}{}\fi
\noindent\textit{Source reference.} Original lemma 0C3P: sources/stacks-chow/tags.txt, source lines 15199--15199. Prerequisite pointer invoked at proof line(s) 4509; it names a standard fact of a sibling chapter that the counted argument uses, and every step of the counted proof built on it is retained verbatim in the delivered body.\par\medskip
\ifdefined\hypertarget\hypertarget{sourcepointer-17}{}\fi
\noindent\textit{Source reference.} Original lemma 0384: sources/stacks-chow/tags.txt, source lines 4223--4223. Prerequisite pointer invoked at proof line(s) 4451; it names a standard fact of a sibling chapter that the counted argument uses, and every step of the counted proof built on it is retained verbatim in the delivered body.\par\medskip
\begin{lemma}
\label{lemma-nodal}
Let $k$ be a field. Let $X$ be a $1$-dimensional algebraic $k$-scheme.
Let $x \in X$ be a closed point. The following are equivalent
\begin{enumerate}
\item $x$ is a node,
\item $k \to \mathcal{O}_{X, x}$ is as in Lemma \ref{lemma-nodal-algebraic},
\item any $\overline{x} \in X_{\overline{k}}$ mapping to $x$ defines
a nodal singularity,
\item $\kappa(x)/k$ is separable, $\mathcal{O}_{X, x}$ is reduced, and
the first Fitting ideal of $\Omega_{X/k}$ generates $\mathfrak m_x$
in $\mathcal{O}_{X, x}$,
\item $\kappa(x)/k$ is separable, $\text{depth}(\mathcal{O}_{X, x}) = 1$, and
the first Fitting ideal of $\Omega_{X/k}$ generates $\mathfrak m_x$
in $\mathcal{O}_{X, x}$,
\item $\kappa(x)/k$ is separable and $\mathcal{O}_{X, x}$ is reduced, has
$\delta$-invariant $1$, and has $2$ geometric branches.
\end{enumerate}
\end{lemma}

\begin{proof}
First assume that $k$ is algebraically closed.
In this case the equivalence of (1) and (3) is trivial.
The equivalence of (1) and (3) with (2) holds because the only
nondegenerate quadric in two variables is $xy$ up to change in
coordinates. The equivalence of (1) and (6) is 
Lemma \ref{lemma-multicross-algebra}.
After replacing $X$ by an affine neighbourhood
of $x$, we may assume there is a closed immersion $X \to \mathbf{A}^n_k$
mapping $x$ to $0$. Let $f_1, \ldots, f_m \in k[x_1, \ldots, x_n]$
be generators for the ideal $I$ of $X$ in $\mathbf{A}^n_k$.
Then $\Omega_{X/k}$ corresponds to the $R = k[x_1, \ldots, x_n]/I$-module
$\Omega_{R/k}$ which has a presentation
$$
R^{\oplus m} \xrightarrow{(\partial f_j/\partial x_i)}
R^{\oplus n} \to \Omega_{R/k} \to 0
$$
(See Algebra, Sections \ref{algebra-section-differentials} and
\ref{algebra-section-netherlander}.)
The first Fitting ideal of $\Omega_{R/k}$ is thus the ideal
generated by the $(n - 1) \times (n - 1)$-minors of the
matrix $(\partial f_j/\partial x_i)$. Hence (2), (4), (5)
are equivalent by Lemma \ref{lemma-fitting-ideal} applied
to the completion of $k[x_1, \ldots, x_n] \to R$
at the maximal ideal $(x_1, \ldots, x_n)$.

\medskip\noindent
Now assume $k$ is an arbitrary field.
In cases (2), (4), (5), (6) the residue field $\kappa(x)$ is
separable over $k$. Let us show this holds as well
in cases (1) and (3). Namely, let $Z \subset X$ be the closed subscheme
of $X$ defined by the first Fitting ideal of $\Omega_{X/k}$.
The formation of $Z$ commutes with field extension
(Divisors, Lemma \ref{divisors-lemma-base-change-and-fitting-ideal-omega}).
If (1) or (3) is true, then there exists a point
$\overline{x}$ of $X_{\overline{k}}$ such that $\overline{x}$
is an isolated point of multiplicity $1$ of $Z_{\overline{k}}$ (as we have the
equivalence of the conditions of the lemma over $\overline{k}$).
In particular $Z_{\overline{x}}$ is geometrically reduced at $\overline{x}$
(because $\overline{k}$ is algebraically closed). Hence
$Z$ is geometrically reduced at $x$
(Varieties, Lemma \ref{varieties-lemma-geometrically-reduced-upstairs}).
In particular, $Z$ is reduced at $x$, hence $Z = \Spec(\kappa(x))$
in a neighbourhood of $x$ and $\kappa(x)$ is geometrically reduced
over $k$. This means that $\kappa(x)/k$ is separable
(Algebra, Lemma \ref{algebra-lemma-characterize-separable-field-extensions}).

\medskip\noindent
The argument of the previous paragraph shows that if (1) or (3) holds, then
the first Fitting ideal of $\Omega_{X/k}$ generates $\mathfrak m_x$.
Since $\mathcal{O}_{X, x} \to \mathcal{O}_{X_{\overline{k}}, \overline{x}}$
is flat and since $\mathcal{O}_{X_{\overline{k}}, \overline{x}}$
is reduced and has depth $1$, we see that (4) and (5) hold
(use Algebra, Lemmas \ref{algebra-lemma-descent-reduced} and
\ref{algebra-lemma-apply-grothendieck}).
Conversely, (4) implies (5) by
Algebra, Lemma \ref{algebra-lemma-criterion-reduced}.
If (5) holds, then $Z$ is geometrically reduced at $x$
(because $\kappa(x)/k$ separable and $Z$ is $x$ in a neighbourhood).
Hence $Z_{\overline{k}}$ is reduced at any point $\overline{x}$
of $X_{\overline{k}}$ lying over $x$. In other words, the first
fitting ideal of $\Omega_{X_{\overline{k}}/\overline{k}}$ generates
$\mathfrak m_{\overline{x}}$ in $\mathcal{O}_{X_{\overline{k}, \overline{x}}}$.
Moreover, since
$\mathcal{O}_{X, x} \to \mathcal{O}_{X_{\overline{k}}, \overline{x}}$ is flat
we see that $\text{depth}(\mathcal{O}_{X_{\overline{k}}, \overline{x}}) = 1$
(see reference above).
Hence (5) holds for $\overline{x} \in X_{\overline{k}}$ and we
conclude that (3) holds (because of the equivalence over algebraically
closed fields). In this way we see that (1), (3), (4), (5)
are equivalent.

\medskip\noindent
The equivalence of (2) and (6) follows from
Lemma \ref{lemma-2-branches-delta-1}.

\medskip\noindent
Finally, we prove the equivalence of (2) $=$ (6) with
(1) $=$ (3) $=$ (4) $=$ (5). First we note that the geometric number
of branches of $X$ at $x$ and the geometric number of branches
of $X_{\overline{k}}$ at $\overline{x}$ are equal by
Varieties, Lemma
\ref{varieties-lemma-geometric-branches-and-change-of-fields}.
We conclude from the information available to us at this point
that in all cases this number is equal to $2$.
On the other hand, in case (1) it is clear that $X$ is geometrically
reduced at $x$, and hence
$$
\delta\text{-invariant of }X\text{ at }x \leq
\delta\text{-invariant of }X_{\overline{k}}\text{ at }\overline{x}
$$
by Varieties, Lemma \ref{varieties-lemma-delta-invariant-and-change-of-fields}.
Since in case (1) the right hand side is $1$, this
forces the $\delta$-invariant of $X$ at $x$ to be $1$
(because if it were zero, then $\mathcal{O}_{X, x}$ would
be a discrete valuation ring by
Varieties, Lemma \ref{varieties-lemma-delta-invariant-is-zero}
which is unibranch, a contradiction). Thus (5) holds.
Conversely, if (2) $=$ (5) is true, then assumptions (a), (b), (c) of
Varieties, Lemma \ref{varieties-lemma-geometrically-normal-in-codim-1}
hold for $x \in X$ by
Lemma \ref{lemma-2-branches-delta-1}. Thus
Varieties, Lemma
\ref{varieties-lemma-delta-invariant-and-change-of-fields-better}
applies and shows that we have equality in the above displayed inequality.
We conclude that (5) holds for $\overline{x} \in X_{\overline{k}}$
and we are back in case (1) by the equivalence of the conditions
over an algebraically closed field.
\end{proof}

\end{document}
